\section{Introduction}
The fundamental challenge in modern memory hierarchies is bridging
the latency gap between processing units and backing stores, a
problem famously characterized as the "Memory Wall"
\cite{Wulf1995Hitting}. To mitigate this, systems employ multi-level
cache hierarchies where the efficiency of the Last Level Cache (LLC)
or software-managed cache dictates overall system performance
\cite{Prieto2011Multilevel}. In these environments, the primary goal
of a replacement policy is to maximize the hit ratio by accurately
identifying and retaining the workload's active "working set"
\cite{Denning2005Locality}.

While Least Recently Used (LRU) algorithms generally provide high hit
ratios, they incur significant synchronization overhead in multi-core
environments due to the need for fine-grained locking on every access
\cite{Herlihy2006The}. Consequently, the CLOCK algorithm and its
concurrent variants, such as NbQ-CLOCK, have become the standard for
high-throughput systems \cite{Eads2014NbQCLOCK}. However, CLOCK
relies on a coarse-grained "blind promotion" mechanism: any accessed
item is given a "second chance" (reference bit set to 1). This
simplicity leaves CLOCK vulnerable to cache pollution caused by scans
or "one-hit wonders"—items accessed once and never again
\cite{Jiang2005CLOCKPro}. When these low-utility items are blindly
promoted, they displace high-utility data, degrading the cache's
effective capacity and overall hit ratio.

Addressing this inefficiency requires distinguishing between
frequency and mere recency. Advanced heuristics like CAR (Clock with
Adaptive Replacement) \cite{Bansal2004CAR} and CLOCK-Pro
\cite{Jiang2005CLOCKPro} attempt to filter these items using history
lists and test segments. However, modern workloads, particularly web
traffic characterized by Zipf-like distributions, exhibit complex
access patterns that static heuristics often fail to capture
optimally \cite{Breslau1999Web}. Recent studies suggest that deciding
what to keep (promotion policy) is often more critical for efficiency
than deciding what to evict \cite{chen2025lazy, Yang2023FIFO}.

In this paper, we introduce PredCLOCK (Predictive CLOCK), a mechanism
to increase the promotion efficiency of the CLOCK algorithm. Rather
than granting a "second chance" to every accessed item, PredCLOCK
utilizes a lightweight Supervised Learning model to predict the reuse
probability of an item at the moment of access. This allows the
system to filter useless promotions, effectively acting as a
gatekeeper that protects the cache from pollution. By leveraging
recent advances in efficient ML inference \cite{Lin2019ONNC}, we
demonstrate that PredCLOCK achieves higher cache utilization
efficiency and robustness against scan-heavy workloads compared to
standard CLOCK implementations.
